\unnsection{Постановка задачи}

    В исторической справке было показано, как уравнение конвекции--реакции--диффузии
    сложилось в единый математический объект, объединивший молекулярную диффузию,
    конвективный перенос и химическую реакцию. В настоящем разделе это уравнение
    формулируется как корректная краевая задача, после чего обзорно рассматриваются
    основные подходы к её решению. Раздел носит \textit{обзорный} характер: его цель ---
    систематизировать способы решения уравнения и указать границы их применимости,
    а не довести до конца вывод какого-либо одного из них.

    \unnsubsection{Математическая постановка краевой задачи}

        Пусть $\Omega \subset \mathbb{R}^n$~--- пространственная область с границей $\partial\Omega$,
        а $(0, T]$~--- рассматриваемый промежуток времени.
        Искомая величина~--- концентрация примеси $c(\mathbf{x}, t)$,
        удовлетворяющая в области $\Omega \times (0, T]$ уравнению
        \[
            \frac{\partial c}{\partial t}
            + \mathbf{u} \cdot \nabla c
            = \nabla \cdot (D\,\nabla c)
            + R(c),
        \]
        где $\mathbf{u}(\mathbf{x}, t)$~--- поле скоростей среды \textit{(конвективный перенос)},
        $D$~--- коэффициент диффузии \textit{(в общем случае тензор)},
        а $R(c)$~--- реакционный член.

        Для однозначной разрешимости уравнение дополняется \textit{начальным условием}
        \[
            c(\mathbf{x}, 0) = c_0(\mathbf{x}), \qquad \mathbf{x} \in \Omega,
        \]
        и \textit{граничными условиями} на $\partial\Omega$, которые традиционно
        разделяют на три типа. Условие \textit{первого рода} \textit{(Дирихле)} задаёт
        концентрацию на границе, $c\big|_{\partial\Omega} = g$; условие \textit{второго рода}
        \textit{(Неймана)} задаёт диффузионный поток,
        $-D\,\dfrac{\partial c}{\partial n}\big|_{\partial\Omega} = q$
        \textit{(в частности, $q = 0$ отвечает непроницаемой стенке)};
        условие \textit{третьего рода} \textit{(Робена)} описывает массообмен
        с окружением, $-D\,\dfrac{\partial c}{\partial n}\big|_{\partial\Omega}
        = \beta\,(c - c_\infty)$.
        Совокупность <<уравнение~$+$~начальное условие~$+$~граничные условия>>
        образует корректно поставленную \textit{(по Адамару)} краевую задачу.

        Для аналитического изучения и для примера расчёта удобна одномерная
        каноническая форма уравнения с постоянными коэффициентами:
        \[
            \frac{\partial c}{\partial t}
            + v\,\frac{\partial c}{\partial x}
            = D\,\frac{\partial^2 c}{\partial x^2}
            + R(c),
            \qquad x \in (0, L), \quad t > 0.
        \]

        Характер задачи и выбор метода решения во многом определяются двумя
        безразмерными параметрами:
        \[
            \mathrm{Pe} = \frac{v L}{D}, \qquad
            \mathrm{Da} = \frac{\tau_{\mathrm{conv}}}{\tau_{\mathrm{reac}}}.
        \]
        Число Пекле $\mathrm{Pe}$ выражает соотношение конвективного и диффузионного
        переноса: при $\mathrm{Pe} \ll 1$ преобладает диффузия и уравнение близко к
        классическому уравнению теплопроводности, а при $\mathrm{Pe} \gg 1$ преобладает
        конвекция и уравнение приобретает почти гиперболический, транспортный характер.
        Число Дамкёлера $\mathrm{Da}$ выражает отношение характерного времени переноса
        к характерному времени реакции, то есть относительную интенсивность
        реакционного слагаемого.

    \unnsubsection{Аналитические методы}

        Решение в замкнутой форме удаётся получить лишь для узкого класса задач:
        линейный \textit{(или отсутствующий)} реакционный член, постоянные коэффициенты,
        простая геометрия области.

        \textbf{Линейное уравнение с постоянными коэффициентами.}
        Для \textit{чистой диффузии} ($v = 0$, $R = 0$) в ограниченной области применяется
        метод разделения переменных \textit{(метод Фурье)}, а в неограниченной --- фундаментальное
        решение в виде ядра теплопроводности. Уравнение \textit{конвекции--диффузии}
        с постоянной скоростью $v$ переходом в подвижную систему координат
        $\xi = x - v t$ сводится к чисто диффузионному; его фундаментальное решение ---
        сносимый потоком гауссов профиль
        \[
            c(x, t) = \frac{1}{\sqrt{4\pi D t}}\,
            \exp\!\left(-\frac{(x - v t)^2}{4 D t}\right).
        \]
        \textit{Линейная реакция} первого порядка $R = -k c$ устраняется подстановкой
        $c = e^{-k t}\,w(x, t)$, после чего задача снова сводится к конвекции--диффузии.
        Общие линейные задачи решаются интегральными преобразованиями
        \textit{(Фурье, Лапласа)} и методом функций Грина~\cite{tihonovSamarskii1999}.

        \textbf{Нелинейная реакция.}
        Ключевой приём здесь~--- поиск решения в виде \textit{бегущей волны}
        $c(x, t) = \varphi(x - \lambda t)$: такая подстановка превращает уравнение
        в частных производных в обыкновенное дифференциальное уравнение для профиля
        $\varphi(\xi)$. Именно так получаются фронты Колмогорова--Петровского--Пискунова---Фишера
        с минимальной скоростью $\lambda_0 = 2\sqrt{D\,\alpha}$, встречавшиеся в
        исторической справке~\cite{kpp1937}, и фронты горения
        Зельдовича--Франк-Каменецкого~\cite{zeldovich1944}. Бегущая волна и есть
        математический образ фронта пламени. Помимо неё используются автомодельные
        \textit{(self-similar)} решения, опирающиеся на инвариантность задачи
        относительно растяжений, а для уравнений с нелинейной конвекцией
        \textit{(типа уравнения Бюргерса)}~--- линеаризующее преобразование Коула--Хопфа.

        Аналитические методы дают точные, физически прозрачные результаты, однако
        охватывают лишь идеализированные постановки. В большинстве практических задач
        приходится прибегать к численному решению.

    \unnsubsection{Численные методы}

        При переменном поле скоростей, нелинейном реакционном члене, сложной геометрии
        области или нескольких пространственных измерениях замкнутого решения не существует,
        и уравнение решается приближённо.

        \textbf{Метод конечных разностей.}
        Область покрывается сеткой по пространству и времени, а производные заменяются
        разностными отношениями. \textit{Явные} схемы \textit{(прямой метод Эйлера)} просты
        в реализации, но условно устойчивы~--- шаг по времени ограничен условием типа
        Куранта. \textit{Неявные} схемы \textit{(обратный метод Эйлера, схема Кранка--Николсон)}
        безусловно устойчивы ценой решения системы линейных алгебраических уравнений
        на каждом шаге. Особую трудность создаёт конвективное слагаемое при больших
        числах Пекле: центральные разности порождают нефизичные осцилляции решения,
        поэтому применяют противопотоковые \textit{(upwind)} схемы и схемы с искусственной
        вязкостью~\cite{samarskii1977}.

        \textbf{Метод прямых.}
        Дискретизация только по пространству превращает уравнение в частных производных
        в систему обыкновенных дифференциальных уравнений по времени, которую решают
        стандартными интеграторами для жёстких систем.

        \textbf{Методы расщепления.}
        Поскольку оператор уравнения естественно распадается на три качественно различные
        части~--- конвекцию, диффузию и реакцию,~--- их удобно учитывать поочерёдно,
        решая каждую подзадачу наиболее подходящим для неё методом
        \textit{(расщепление Марчука--Стрэнга)}~\cite{marchukSplitting1988}. Этот приём
        особенно органичен именно для уравнения конвекции--реакции--диффузии.

        \textbf{Методы конечных элементов и конечных объёмов.}
        Метод конечных элементов удобен для областей сложной геометрии, а метод конечных
        объёмов обеспечивает точное выполнение закона сохранения вещества, что важно
        для задач переноса.

    \unnsubsection{Феноменологические и экспериментально-статистические подходы}

        В наиболее сложных режимах фронт реакции теряет устойчивость и фрагментируется
        \textit{(интенсивная реакция, неустойчивости Дарье--Ландау и диффузионно-тепловая
        неустойчивость при малых числах Льюиса)}. Решение задачи становится сложной,
        стохастической, фрактально-подобной структурой, и его прямое вычисление теряет
        практический смысл. В этом случае применяют иную стратегию: не решают уравнение
        непосредственно, а описывают \textit{макроскопическое поведение} его решения
        эмпирическими скейлинговыми соотношениями и статистическим \textit{(фрактальным)}
        анализом экспериментальных или расчётных данных.

        Характерный пример такого подхода~--- работа П. В. Москалёва, В. П. Денисенко
        и И. А. Кириллова~\cite{moskalev2023}, посвящённая динамике ультрабедных
        водородо-воздушных пламён в горизонтальной цилиндрической ячейке Хеле--Шоу.
        Горение представляет собой полную <<триаду>> процессов уравнения:
        реакцию \textit{(сгорание по закону Аррениуса)}, диффузию \textit{(с малым числом Льюиса)}
        и конвекцию \textit{(плавучесть, для компенсации которой ячейку располагают
        горизонтально)}. В простейшем случае фронт пламени~--- это бегущая волна, то есть
        решение уравнения; однако в ультрабедном режиме фронт распадается на дискретные
        <<дрейфующие шаровые пламёна>> и образует фрактальную структуру, закономерно
        меняющуюся с ростом концентрации водорода \textit{(рис.~\ref{fig:morphotypes})}.

        \begin{figure}[H]
            \centering
            \includegraphics[width=\linewidth]{morphotypes}
            \caption{Три характерных морфотипа следов ультрабедного водородо-воздушного
                пламени в горизонтальной цилиндрической ячейке Хеле--Шоу:
                \textit{а)}~лучеобразные \textit{(6,8\,об.\,\% водорода)};
                \textit{б)}~дендритные \textit{(7,2\,об.\,\%)};
                \textit{в)}~квазиоднородные \textit{(9,0\,об.\,\%)}.
                Источник: \cite{moskalev2023}}
            \label{fig:morphotypes}
        \end{figure}

        Поэтому авторы характеризуют динамику не решением уравнения, а набором
        аппроксимирующих соотношений, оцениваемых по экспериментальным данным методом
        наименьших квадратов: фрактальной размерностью следа $d_B(x)$
        \textit{(метод клеточного покрытия)}, аппроксимированной логистической моделью;
        степенной и экспоненциальной моделями длины пути фронта $s(t)$; и
        логарифмически-экспоненциальной моделью зависимости начальной скорости фронта
        $v_0$ от объёмной концентрации водорода $x$ \textit{(рис.~\ref{fig:velocity})}.
        Эти скейлинговые соотношения играют роль эффективного <<решения>> задачи там,
        где уравнение не поддаётся прямому интегрированию.

        \begin{figure}[H]
            \centering
            \includegraphics[width=0.68\linewidth]{velocity_v0}
            \caption{Зависимость начальной скорости фронта $v_0$ ультрабедного
                водородо-воздушного пламени от объёмной концентрации водорода $x$
                и её логарифмически-экспоненциальная аппроксимация. Источник:
                \cite{moskalev2023}}
            \label{fig:velocity}
        \end{figure}

        Таким образом, спектр методов простирается от точных аналитических решений
        \textit{(узкий класс идеализированных задач)} через универсальные численные схемы
        к феноменологическому описанию, применяемому в наиболее сложных режимах. Выбор
        метода определяется параметрами задачи \textit{($\mathrm{Pe}$, $\mathrm{Da}$)},
        геометрией области и требуемой точностью. В следующем разделе один из подходов
        доводится до конкретного примера расчёта с приведением кода.
